\documentclass[11pt,a4paper]{article}
\usepackage{fontspec,amsmath,amssymb,amsthm,microtype}
\setmainfont{lmroman10-regular.otf}[BoldFont=lmroman10-bold.otf,ItalicFont=lmroman10-italic.otf,BoldItalicFont=lmroman10-bolditalic.otf]
\usepackage{xcolor}
\usepackage[most]{tcolorbox}
\usepackage{mathrsfs}
\providecommand{\lt}{<}
\providecommand{\gt}{>}
\setmonofont{DejaVuSansMono.ttf}[BoldFont=DejaVuSansMono-Bold.ttf,ItalicFont=DejaVuSansMono-Oblique.ttf,BoldItalicFont=DejaVuSansMono-BoldOblique.ttf,Scale=0.86]
\newfontfamily\leanfallback{STIXTwoMath-Regular.otf}[Scale=0.86]
\definecolor{leanrule}{gray}{0.62}
\definecolor{leanback}{gray}{0.962}
\definecolor{leantitle}{gray}{0.22}
\definecolor{leanmuted}{gray}{0.38}
\newtcolorbox{leanbox}[1]{enhanced,breakable,colback=leanback,colframe=leanrule,
  boxrule=0.4pt,arc=0.6mm,left=2.4mm,right=2.4mm,top=1.8mm,bottom=1.8mm,
  before skip=0.8em,after skip=1em,fonttitle=\small\bfseries,coltitle=leantitle,
  colbacktitle=leanback,titlerule=0.3pt,title={#1}}
\newlength{\leancharwidth}
\newenvironment{leancode}{\par\smallskip\begingroup\ttfamily\footnotesize\raggedright
  \settowidth{\leancharwidth}{M}\setlength{\parskip}{0pt}\setlength{\parindent}{0pt}
  \frenchspacing}{\par\endgroup\smallskip}
\newcommand{\leanline}[2]{\hangindent=\dimexpr#1\leancharwidth+4\leancharwidth\relax
  \hangafter=1\noindent\hspace*{#1\leancharwidth}#2\par}
\newcommand{\leansource}[1]{\par{\footnotesize\color{leanmuted}#1}\par}
\newcommand{\leanlabel}[1]{\par\smallskip\noindent{\small\bfseries #1}\par}
\usepackage[margin=25mm]{geometry}
\usepackage[unicode,hidelinks]{hyperref}
\hypersetup{pdfauthor={LangevinLean Project}}
\setlength{\emergencystretch}{3em}
\setlength{\parskip}{0.4em}
\setlength{\parindent}{0pt}
\allowdisplaybreaks[2]
\raggedbottom
\title{\texorpdfstring{The explicit actual Gaussian Metropolis kernel}{The explicit actual Gaussian Metropolis kernel}}
\author{LangevinLean Project}
\date{}
\begin{document}\maketitle
Draft mathematical exposition; these notes may contain mistakes.


The transition kernel samples an actual fresh standard Gaussian, maps accepted mass to the Langevin proposal and rejected mass to the current state. Its action on every initial law equals the existing Gaussian MALA law step exactly.


\section{\texorpdfstring{Constructing the actual fresh-noise transition}{Constructing the actual fresh-noise transition}}


\textbf{Construction.} At state \(x\), take the product \(\delta_x\otimes\gamma\), where \(\gamma\) is standard Gaussian. Weight by actual Metropolis acceptance and rejection probabilities and map to the proposed or unchanged state. The resulting explicit kernel at \(x\) is exactly the MALA law step applied to \(\delta_x\).



\textbf{Proof.} The product kernel is Markov. Measurability of the actual gradient proposal and acceptance permits weighted mapping. Evaluating at a state and substituting the Dirac initial law yields the same accepted/\allowbreak{}rejected formula as the actual MALA step.

\begin{leanbox}{Lean: \texttt{gaussianPairKernel}, \texttt{gaussianPairKernel\_\allowbreak{}isMarkov}, \texttt{explicitKernel} and 1 more}
\leansource{\texttt{gaussianPairKernel} \quad WassersteinMixing/\allowbreak{}SemigroupMALAKernel.lean, lines 24–25}
\begin{leancode}
\leanline{0}{def\ gaussianPairKernel\ :\ Kernel\ E\ (E\ ×\ E)\ :=}
\leanline{2}{Kernel.id\ ×ₖ\ Kernel.const\ E\ (stdGaussian\ E)}
\end{leancode}
\leansource{\texttt{gaussianPairKernel\_\allowbreak{}isMarkov} \quad WassersteinMixing/\allowbreak{}SemigroupMALAKernel.lean, lines 27–29}
\begin{leancode}
\leanline{0}{instance\ gaussianPairKernel\_isMarkov\ :\ IsMarkovKernel\ (gaussianPairKernel\ (E\ :=\ E))\ :=\ by}
\leanline{2}{unfold\ gaussianPairKernel}
\leanline{2}{infer\_instance}
\end{leancode}
\leansource{\texttt{explicitKernel} \quad WassersteinMixing/\allowbreak{}SemigroupMALAKernel.lean, lines 32–37}
\begin{leancode}
\leanline{0}{def\ explicitKernel\ (V\ :\ E\ →\ ℝ)\ (h\ :\ ℝ)\ :\ Kernel\ E\ E\ :=}
\leanline{2}{((gaussianPairKernel\ (E\ :=\ E)).withDensity}
\leanline{6}{(fun\ \_\ p\ ↦\ ENNReal.ofReal\ (ULA.malaProposalAcceptance\ V\ h\ p))).map}
\leanline{8}{(fun\ p\ ↦\ ULA.malaProposal\ V\ h\ p.1\ p.2)\ +}
\leanline{2}{((gaussianPairKernel\ (E\ :=\ E)).withDensity}
\leanline{6}{(fun\ \_\ p\ ↦\ ENNReal.ofReal\ (1\ -\ ULA.malaProposalAcceptance\ V\ h\ p))).map\ Prod.fst}
\end{leancode}
\leansource{\texttt{explicitKernel\_\allowbreak{}apply} \quad WassersteinMixing/\allowbreak{}SemigroupMALAKernel.lean, lines 40–54}
\begin{leancode}
\leanline{0}{theorem\ explicitKernel\_apply\ (V\ :\ E\ →\ ℝ)\ (h\ :\ ℝ)\ (L\ :\ ℝ≥0)}
\leanline{4}{(hdiff\ :\ Differentiable\ ℝ\ V)\ (hLip\ :\ LipschitzWith\ L\ (gradient\ V))\ (x\ :\ E)\ :}
\leanline{4}{explicitKernel\ V\ h\ x\ =\ ULA.malaOneStep\ V\ h\ (Measure.dirac\ x)}
\end{leancode}
\leanlabel{How the Lean proof follows the argument}
\texttt{gaussian\allowbreak{}Pair\allowbreak{}Kernel} is the product of the identity and constant-Gaussian kernels. \texttt{explicit\allowbreak{}Kernel} is the sum of its actual weighted mapped kernels. \texttt{explicit\allowbreak{}Kernel\_\allowbreak{}apply} compares their event masses with the Dirac-start MALA step.
\end{leanbox}


Earlier results used here: \cite[Section~3.34.1]{LangevinLean2026ULASC}; \cite[Section~3.36.2]{LangevinLean2026ULASC}; \cite[Section~3.36.5]{LangevinLean2026ULASC}; \cite[Section~3.36.6]{LangevinLean2026ULASC}; \cite[Section~3.36.8]{LangevinLean2026ULASC}.

\section{\texorpdfstring{The exact law action and mass normalization}{The exact law action and mass normalization}}


\textbf{Claim.} For every sigma-finite initial law \(\mu\), the explicit kernel action is the actual MALA output law. It has total mass one at every state.



\textbf{Proof.} Composing the fresh-Gaussian pair kernel with \(\mu\) gives \(\mu\otimes\gamma\). Mapping and weighting commute with actual measure/\allowbreak{}kernel composition, so the explicit output is exactly the accepted/\allowbreak{}rejected law formula for \(\mu\). For a Dirac input, acceptance plus rejection partition the entire Gaussian mass, proving the Markov normalization.

\begin{leanbox}{Lean: \texttt{gaussianPairKernel\_\allowbreak{}comp}, \texttt{explicitKernel\_\allowbreak{}comp}, \texttt{explicitKernel\_\allowbreak{}isMarkov}}
\leansource{\texttt{gaussianPairKernel\_\allowbreak{}comp} \quad WassersteinMixing/\allowbreak{}SemigroupMALAKernel.lean, lines 58–60}
\begin{leancode}
\leanline{0}{theorem\ gaussianPairKernel\_comp\ (μ\ :\ Measure\ E)\ [SFinite\ μ]\ :}
\leanline{4}{gaussianPairKernel\ (E\ :=\ E)\ ∘ₘ\ μ\ =\ μ.prod\ (stdGaussian\ E)}
\end{leancode}
\leansource{\texttt{explicitKernel\_\allowbreak{}comp} \quad WassersteinMixing/\allowbreak{}SemigroupMALAKernel.lean, lines 64–77}
\begin{leancode}
\leanline{0}{theorem\ explicitKernel\_comp\ (V\ :\ E\ →\ ℝ)\ (h\ :\ ℝ)\ (L\ :\ ℝ≥0)}
\leanline{4}{(hdiff\ :\ Differentiable\ ℝ\ V)\ (hLip\ :\ LipschitzWith\ L\ (gradient\ V))}
\leanline{4}{(μ\ :\ Measure\ E)\ [SFinite\ μ]\ :}
\leanline{4}{explicitKernel\ V\ h\ ∘ₘ\ μ\ =\ ULA.malaOneStep\ V\ h\ μ}
\end{leancode}
\leansource{\texttt{explicitKernel\_\allowbreak{}isMarkov} \quad WassersteinMixing/\allowbreak{}SemigroupMALAKernel.lean, lines 80–98}
\begin{leancode}
\leanline{0}{theorem\ explicitKernel\_isMarkov\ (V\ :\ E\ →\ ℝ)\ (h\ :\ ℝ)\ (L\ :\ ℝ≥0)}
\leanline{4}{(hdiff\ :\ Differentiable\ ℝ\ V)\ (hLip\ :\ LipschitzWith\ L\ (gradient\ V))\ :}
\leanline{4}{IsMarkovKernel\ (explicitKernel\ V\ h)}
\end{leancode}
\leanlabel{How the Lean proof follows the argument}
\texttt{gaussian\allowbreak{}Pair\allowbreak{}Kernel\_\allowbreak{}comp} identifies the actual independent-product sampling law. \texttt{explicit\allowbreak{}Kernel\_\allowbreak{}comp} uses map, density and sum composition identities. \texttt{explicit\allowbreak{}Kernel\_\allowbreak{}is\allowbreak{}Markov} proves unit mass from the true accepted/\allowbreak{}rejected partition.
\end{leanbox}


Earlier results used here: \cite[Section~3.34.1]{LangevinLean2026ULASC}; \cite[Section~3.35.1]{LangevinLean2026ULASC}; \cite[Section~3.35.2]{LangevinLean2026ULASC}; \cite[Section~3.35.4]{LangevinLean2026ULASC}; \cite[Section~3.36.2]{LangevinLean2026ULASC}; \cite[Section~3.36.3]{LangevinLean2026ULASC}; \cite[Section~3.36.4]{LangevinLean2026ULASC}; \cite[Section~3.36.5]{LangevinLean2026ULASC}; \cite[Section~3.36.6]{LangevinLean2026ULASC}; \cite[Section~3.36.8]{LangevinLean2026ULASC}.

\begin{thebibliography}{ULA-LSI}
\bibitem[ULA-SC]{LangevinLean2026ULASC}
LangevinLean Project. \emph{Unadjusted Langevin under Strong Convexity}. LL-01, 2026. Draft report, content revision d34ead06987c. \url{https://odakyildiz.github.io/LangevinLean/docs/pdf/UnadjustedLangevin/FullProof.pdf}.
\end{thebibliography}

\end{document}
